\section{Introduction and main results}\label{sec:introduction}

In a jointly $C^2$ unfolding, a periodic orbit whose return displacement
has an ordinary double zero has at most two nearby cycles: the second
derivative keeps its sign, and Rolle's theorem applies~\cite{roussarie1998}.
We study a continuous piecewise smooth planar system in which the reference
displacement has such a double zero, but the sharp local cyclicity is six.

The reference orbit has two spatially separated quadratic grazings, with both
active sides agreeing after transport to one transverse section. Each
passage contributes a term of order $q^{3/2}$ in a transported transverse
coordinate $q$~\cite{dibbc2001,nordmark2002}, and their nonzero leading
coefficients cancel. The quadratic defect of the composed passages cancels
as well, as the expansion below shows. The reference displacement is therefore $C^2$, with
its nonzero quadratic term supplied by the smooth closing map. Under
perturbation, the strength sum can become nonzero and the activation
thresholds can separate, so the family need not be jointly $C^2$.
Together with the closing offset and multiplier defect, these give four
scalar return-map controls. We answer three questions.
\begin{enumerate}
\item[(Q1)] What is the cyclicity of the reference cycle?
\item[(Q2)] What parameter scales are necessary for six limit cycles,
and what is the limiting shape of the six-cycle parameter set?
\item[(Q3)] Does a fourfold root of the limiting polynomial continue
to a cycle of multiplicity four, and what is its local unfolding?
\end{enumerate}
Theorems~\ref{thm:six-main}, \ref{thm:main-saturation-summary}
and~\ref{thm:actual-quartic-summary} below answer these questions.

\subsection{The return map}

In each of the two grazing boxes, the planar equation has the local form
\[
 \dot{\mathbf{x}}=X^-_\lambda(\mathbf{x})
 +(h_\lambda(\mathbf{x}))_+Z_\lambda(\mathbf{x}),
 \qquad \mathbf{x}\in\R^2,\qquad u_+=\max\{u,0\}.
\]
The local data are jointly smooth in state and parameters and may differ
between the boxes. The active branch
$X^+_\lambda=X^-_\lambda+h_\lambda Z_\lambda$ agrees with $X^-_\lambda$
on the seam $\{h_\lambda=0\}$, so the field is continuous there.
The boxes are joined by regular flow arcs; Definition~\ref{def:admissible-family}
gives the precise admissibility conditions. The threshold-harvesting family
of Theorem~\ref{thm:harvesting-realization} realizes these conditions.

In a transverse coordinate $q$ constructed in Section~\ref{sec:passage}, the
return map and the displacement have the form
\[
 P(q)=R\{\delta+\Gamma_2(\Gamma_1(q)-\delta)\},
 \qquad D(q)=P(q)-q.
\]
Here $\Gamma_1,\Gamma_2$ are the passage maps near the two tangencies, $R$ is
the smooth transition along the rest of the cycle, and parameter dependence is
suppressed. The second passage is activated by the intermediate input
$\Gamma_1(q)-\delta$, so the two grazing thresholds on the section are $0$ and
$\Gamma_1^{-1}(\delta)$. Section~\ref{sec:passage} derives this composition
from the vector field and shows that it is the first return on a fixed
interval.

We use four return-map controls,
\[
 a=R(0),\qquad b=R'(0)-1,\qquad m=A_1+A_2,\qquad\delta,
\]
namely the offset, the multiplier defect, the strength sum and the threshold
separation. They are smooth functions of the physical parameters and need not
be independent. All four vanish at the reference parameter, where $A_1=-A_2\ne0$
and $c_0=R''(0)/2\ne0$. Writing $q_+=\max\{q,0\}$, the correction
$\Gamma_\nu(q)-q$ begins with $A_\nu q_+^{3/2}$, and the reversibility of the
quadratic passage fixes its next coefficient at $3A_\nu^2/4$
(Proposition~\ref{prop:quadratic-passage}). Consequently, when the thresholds
are aligned,
\[
 \Gamma_2(\Gamma_1(q))-q
 =mq_+^{3/2}+\tfrac34m^2q_+^2+O(q_+^{5/2})
 \qquad(\delta=0).
\]
Both displayed terms vanish when $m=0$, so the quadratic term of the reference
displacement comes from $R$ (Lemma~\ref{lem:passage:finite-inverse}). For
$\delta\ne0$, the two $q^{3/2}$ contributions start at different thresholds.
Their singular curvatures have opposite signs and compete. Neither an outer
expansion nor the number of active intervals controls the fixed-point count
on the whole return interval; estimates relative to the moving thresholds
are needed.

We use the standard notion of cyclicity~\cite{roussarie1998}.

\begin{definition}[Cyclicity]\label{def:cyclicity}
Let $(X_\lambda)$ be a family of planar vector fields and let $\mathcal C$ be a
periodic orbit of $X_{\lambda_0}$. Let $N(\lambda,\varepsilon)$ be the number
of limit cycles of $X_\lambda$ at Hausdorff distance less than $\varepsilon$
from $\mathcal C$. The cyclicity of $\mathcal C$ in $(X_\lambda)$ is
\[
 \operatorname{Cycl}\bigl((X_\lambda),\mathcal C\bigr)
=\inf_{\varepsilon,\eta>0}\ \sup_{|\lambda-\lambda_0|<\eta}N(\lambda,\varepsilon).
\]
For a class $\mathfrak F$ of families through
$X_{\lambda_0}$, $\operatorname{Cycl}(\mathfrak F,\mathcal C)$ denotes the supremum
over the families in $\mathfrak F$.
\end{definition}

The families we consider are the admissible families of
Definition~\ref{def:admissible-family}: finite-dimensional smooth families of
continuously matched fields that preserve regular switching curves, the two
quadratic tangencies, their transported active side and the one-turn itinerary
of $\mathcal C$, with
$a=b=\delta=0$, $A_1=-A_2\ne0$ and $c_0\ne0$ at the reference parameter.
When the map $\lambda\mapsto(a,b,m,\delta)$ is a submersion at
$\lambda_0$, we fix a local inverse
\[
 \Lambda(a,b,m,\delta)=\lambda(a,b,m,\delta;\xi_0),
\]
with the remaining parameters fixed at $\xi_0$, and call the resulting
fixed four-parameter family the \emph{standard unfolding}.

\subsection{Sharp cyclicity}

\begin{theorem}[Cyclicity]\label{thm:six-main}
Let $\mathfrak F$ be the class of admissible families through the reference
field.
\begin{enumerate}
\item For every family in $\mathfrak F$ there are a neighbourhood $U$ of
$\lambda_0$ and an interval $I=(-q_0,q_0)$ on the transverse section such
that, for every $\lambda\in U$, the return map $P_\lambda$ has at most six
fixed points in $I$; at most four if $m(\lambda)=0$ or $\delta(\lambda)=0$;
and at most two if $m(\lambda)=\delta(\lambda)=0$.
\item If $\lambda\mapsto(a,b,m,\delta)$ is a submersion at $\lambda_0$,
all four bounds are attained by hyperbolic fixed points along smooth
parameter paths in the corresponding slices that tend to $\lambda_0$.
For every sufficiently small $m$ with $mc_0<0$, the standard unfolding
has six hyperbolic fixed points of order $m^2$ in $I$ whenever
$(a/m^4,b/m^2,\delta/m^3)$ belongs to a fixed nonempty open set.
Three cycles are attracting and three repelling.
\item Consequently $\operatorname{Cycl}(\mathfrak F,\mathcal C)=6$. The
corresponding suprema over admissible families contained in $\{m=0\}$, in
$\{\delta=0\}$ and in $\{m=\delta=0\}$ are four, four and two.
\end{enumerate}
\end{theorem}

The proof controls curvature before counting fixed points. With the other
controls fixed, the outer equation $P''(q)=0$ is represented exactly by a
separation graph $\delta=\Delta(r)$, $r=\sqrt q$. This graph has at most one
positive critical point, and hence at most two intersections with a fixed
separation level. Estimates relative to the moving thresholds exclude
further curvature zeros outside the outer region. The two possible
curvature zeros and the two thresholds divide the continuous derivative
$P'$ into at most five strictly monotone pieces. Thus $P'-1$ has at most
five zeros, and Rolle's theorem gives at most six fixed points, including
those at the thresholds.

Theorem~\ref{thm:six-main} is proved in Section~\ref{sec:curvature-graph}
(Proposition~\ref{prop:common-count}, Corollary~\ref{cor:four-slices} and
Propositions~\ref{prop:single-scale-six}--\ref{prop:sharp-slice-attainment})
and Section~\ref{subsec:specified-field}. Proposition~\ref{prop:specified-rank}
constructs, through every reference field, an admissible family for which
$\lambda\mapsto(a,b,m,\delta)$ is a submersion; Corollary~\ref{cor:specified-cyclicity}
gives the same cyclicity in the smaller class of local perturbations used there.

\subsection{The saturation domain}

For an admissible family we call the set of parameters $\lambda\in U$ for
which $P_\lambda$ has six fixed points in $I$ the \emph{saturation domain},
and a sequence of such parameters tending to $\lambda_0$ a \emph{saturated
sequence}. Describing this set requires estimates that remain valid when the
parameters tend to the reference at arbitrary relative rates, including
configurations whose fixed points approach the thresholds. Six zeros of $D$
force five alternating extrema, and their locations and values constrain the
parameters.

Consider the sign case $A_1(\lambda_0)=B>0$, $c_0=-C<0$. Saturation forces
$m>0$, $\delta>0$ and $b<0$. Beyond both thresholds, the balance of
$mq^{3/2}$ and $-Cq^2$ suggests $q$ of order $m^2$. At this scale, $bq$, the
separation term $\tfrac32B\delta\sqrt q$ and the offset $a$ have the same order
when $b,\delta,a$ have weights $2,3,4$ relative to $m$.
Theorem~\ref{thm:main-saturation-summary} shows that every six-cycle
configuration obeys upper bounds with these weights.

Define the normalized parameters and the normalized positive state by
\[
 A_4=\frac{C^3a}{m^4},\qquad S_4=-\frac{Cb}{m^2},\qquad
 T_4=\frac{3BC^2\delta}{2m^3},\qquad x=\frac{C\sqrt q}{m}.
\]
The limiting polynomial is
\[
 \Phi_{A,S,T}(x)=A-Sx^2+x^3-x^4+Tx.
\]
With $e_j$ denoting the elementary symmetric polynomials of
$r=(r_1,r_2,r_3,r_4)$, set
\[
 \mathcal K_4=
 \left\{(-e_4(r),e_2(r),e_3(r)):
 r_i\ge0,\ \sum_{i=1}^4r_i=1\right\}.
\]
This is the compact set of coefficients for which $\Phi_{A,S,T}$ has four
nonnegative real roots, counted with multiplicity. For the standard unfolding
and $m>0$, let $\mathscr T_m$ be the set of $(A,S,T)$ in a fixed compact box
$\mathcal B$ for which the parameters
$(a,b,m,\delta)=(m^4A/C^3,\,-m^2S/C,\,m,\,2m^3T/(3BC^2))$ give six fixed points
in $I$; the box $\mathcal B$ contains all normalized six-cycle parameters
(see~\eqref{def:actual-main-domain}). We write $d_H$ for the Hausdorff distance
in $\R^3$.

\begin{theorem}[Saturation at the principal scale]\label{thm:main-saturation-summary}
\label{thm:k4-domain}
\begin{enumerate}
\item For every admissible family, after shrinking $U$ and $I$ if necessary,
there is $K>0$ such that, whenever
$\lambda\in U$ and $P_\lambda$ has six fixed points $q_1,\dots,q_6$ in $I$,
\[
 0<|\delta|\le K|m|^3,\quad |b|\le K|m|^2,
 \quad |a|\le K|m|^4,\quad \max_j|q_j|\le K|m|^2.
\]
\item For every admissible family in the sign case $A_1(\lambda_0)=B>0$,
$c_0=-C<0$, every accumulation point of $(A_4,S_4,T_4)$ along a saturated
sequence belongs to $\mathcal K_4$.
\item In the same sign case, for the standard unfolding,
\[
 d_H(\overline{\mathscr T_m},\mathcal K_4)\longrightarrow0\quad(m\to0^+).
\]
For every compact set $E\subset\operatorname{int}\mathcal K_4$ there is
$m_E>0$ such that $E\subset\mathscr T_m$ for $0<m<m_E$, and the six fixed
points are simple.
\end{enumerate}
\end{theorem}

For $0<m<m_E$, a point of $E$ gives four fixed points beyond both thresholds
and two with $q<0$, where both passages are inactive. Boundary points of
$\mathcal K_4$ correspond to zero or repeated roots of $\Phi_{A,S,T}$.
Theorem~\ref{thm:actual-quartic-summary} continues the fourfold limiting root
to a nearby parameter branch with one cycle of multiplicity four and two
simple cycles. Theorem~\ref{thm:main-saturation-summary} is proved in
Section~\ref{sec:main-saturation}: part~1 is
Proposition~\ref{sat:weighted-confinement}, part~2 is
Proposition~\ref{prop:k4-necessary}, and part~3 is completed at the end of
Section~\ref{subsec:four-root-limit}.

\subsection{Limit cycles of multiplicity four}

The boundary point
\[
 (A,S,T)=\left(-\frac1{256},\frac38,\frac1{16}\right)
\]
has four roots equal to $1/4$, so $\Phi_{A,S,T}(x)=-(x-1/4)^4$. The next
theorem continues this configuration to the vector field for every small
positive $m$ and counts all nearby limit cycles.

\begin{theorem}[Limit cycles of multiplicity four]\label{thm:actual-quartic-summary}
Consider the standard unfolding in the sign case $A_1(\lambda_0)=B>0$,
$c_0=-C<0$. For every sufficiently small $m>0$ there are parameters
$(a_*(m),b_*(m),\delta_*(m))$ at which $P_\lambda$ has a fixed point of
multiplicity four beyond both thresholds and two simple fixed points with
$q<0$. Near these parameters, with $m$ fixed, smooth changes of the state and
of $(a,b,\delta)$ reduce the fixed-point equation near the fourfold fixed
point, up to a nonvanishing factor, to
\[
 y^4+uy^2+vy+w=0.
\]
Off the set where this quartic has a multiple root, $P_\lambda$ has two, four
or six fixed points in $I$, and each count occurs on a nonempty open set of
parameters. At the fourfold point $P_\lambda$ has exactly three fixed points
in $I$.
\end{theorem}

The proof first locates all zeros of $D$ in $I$
(Lemma~\ref{mainloc:all-roots}): two simple zeros with $q<0$ and a cluster
beyond both thresholds, at distance of order $m^2$ from them. On this cluster
$D$ is smooth in $q$, so the Malgrange preparation theorem applies once the
fourfold zero has been constructed in normalized variables
(Theorem~\ref{thm:actual-quartic}). The exact identities $q_*=r_{\rm c}^2$
and $\delta_*=\delta_{\rm c}$ in that theorem place the constructed cycle at
the critical point of the separation graph. The fold surface, the curve of triple roots
and the curve of two double roots are then explicit
(Proposition~\ref{quartic:strata}). Theorem~\ref{thm:actual-quartic-summary}
follows from Lemma~\ref{mainloc:all-roots} and Theorem~\ref{thm:actual-quartic}.

\subsection{Realization and further results}

Section~\ref{sec:sharpness} shows that the hypotheses hold in two ways: for
local perturbations of an arbitrary reference field
(Proposition~\ref{prop:specified-rank}), and for the natural parameters of a
predator--prey model with continuous threshold harvesting
(Theorem~\ref{thm:harvesting-realization}). The model combines the smooth core
of Huang, Gong and Ruan~\cite{huang2013} with continuous threshold
policies~\cite{bohn2011,leardrebaza2011,rebaza2012,onana2019}. In this model
the reference cycle, the cancellation, the negative closing curvature and the
submersion condition are all obtained with positive ecological parameters.

Online Resource~1 treats the corner $S_4\to0$ of $\mathcal K_4$, where three
roots tend to zero, at a secondary scale. It identifies the limiting cubic
coefficient region, constructs six limit cycles with positive offset $a>0$,
and unfolds a cusp at the triple root. The proofs in the present paper do not
use these results.

\subsection{Relation to previous work}\label{sec:boundaries}

The $q^{3/2}$ correction for continuous fields and the composition of
successive grazing maps belong to the established theory of grazing
bifurcations~\cite{nordmark1991,dibbc2001,nordmark2002,dibernardo2008,colombo2012}. For a
periodically forced oscillator with a continuous piecewise-linear restoring
force, Liu and Yue compose two such passages with regular transitions and test
the second activation at the intermediate state~\cite{liuyue2023}. For
bilateral impacts, Li, Wei, Tan and Xi compose two local impact maps and give a
necessary existence condition~\cite[Theorem~1]{liweitanxi2014}; related
compositions occur in~\cite{lilimeng2021,miao2023}. Two-parameter grazing
bifurcations, including simultaneous contacts at separated points, are
classified in~\cite{kowalczyk2006}. We take the composition rule as known and
determine the cyclicity of the resulting return map.

Translated-power compositions, studied by Dumortier and
Roussarie~\cite{dumortierroussarie2008}, provide a related approach to
counting return-map zeros. Panazzolo's generalized derivation--division
algorithm gives root bounds for these compositions and parameter-analytic
compensator expansions~\cite[arXiv version, Theorem~2.2, Sections~3--4 and
Proposition~5.2]{panazzolo2024}. That explicit function class has not been
matched here to the actual return, with its smooth profiles, remainders and
closing map, across moving activation thresholds. We count its curvature
zeros through the exact separation graph and uniform threshold estimates.

Bifurcations of periodic orbits in planar piecewise smooth systems have been
classified in low codimension~\cite{kuznetsov2003,guardia2011}, and limit
cycles of continuous piecewise linear systems are studied
in~\cite{freire1998}; see also~\cite{simpson2010}. A related configuration is
a periodic orbit tangent to two switching lines in a discontinuous
system with sliding, whose bifurcations resemble those of double and triple
cycles~\cite{roitenberg2025}. Here the field is continuous, there is no
sliding, and we prove cyclicity six for a semistable cycle with cancelling
leading passage coefficients.

Nonhyperbolic cycles and cancellations also occur in other grazing problems.
For one continuous planar grazing contact, Simpson and Meiss obtain a return
with a $q^{3/2}$ term, a nearby fold and a sixth-order separation of the fold
and grazing curves~\cite[Theorem~1]{simpsonmeiss2008}. Colombo and
Dercole~\cite{colombodercole2010} study discontinuity-induced bifurcations of
nonhyperbolic cycles, including a fold with multiplier one. In resonant impact
grazing, a singular transport contribution is suppressed while the local
square-root correction remains nonzero~\cite[Theorems~2.4--2.5]{simpsonghosh2026}.
T.~Matsuhisa (\emph{Induced-contact stability and activation thresholds near
double compliant grazing}, preprint dated 3~October 2026,
\href{https://doi.org/10.2139/ssrn.7556610}{SSRN:7556610}) studies
height-resolved corrections and paths cancelling the direct second-gap height
against the transmitted first-contact height. The continuation of the matched
periodic branch uses the assumption that $D_zP_f-I$ is invertible (Theorem~2).
Here two nonzero transported passage strengths cancel at a reference cycle
with multiplier one, so this continuation argument is unavailable.

Large numbers of limit cycles at discontinuous tangencies arise from a
different local geometry. Novaes, Teixeira and Zeli study a cycle through a
visible two-fold point~\cite{novaesteixeirazeli2018}. Novaes and Silva
construct five hyperbolic limit cycles near one monodromic tangential
singularity of a quadratic Filippov family, with a sixth created by a
pseudo-Hopf perturbation~\cite[Example~3 and Remark~1]{novaessilva2021}; further constructions appear
in~\cite{novaessilva2025}. These tangencies occur at a single point of a
Filippov system, whereas the two contacts here are separated and the fields
are continuously matched. For Sotomayor--Teixeira regularization,
Bonet-Rev\'es and M-Seara treat an identity-return centre and hyperbolic
grazing cycles~\cite[Theorems~2.3--2.4]{bonetseara2016}, and prove a unique
saddle-node for repelling hyperbolic grazing
\cite[Theorems~2.4 and~2.10]{bonetseara2022}.

Two preprints of the author treat related nonhyperbolic returns: an
invariant-line entry--exit followed by a separated quadratic grazing
(H.~Lu, \emph{Neutral Entry--Exit Cycles with Quadratic Grazing: Uniform
Return Reduction and Local Two-Parameter Bifurcations}, preprint, 2026,
\href{https://arxiv.org/abs/2607.27464v4}{arXiv:2607.27464v4}), and a single
high-order grazing with multiwell passages (H.~Lu, \emph{Neutral Returns at
High-Order Grazing: Sharp Cyclicity, Physical Codimension, and Weighted
Crossover}, preprint, 2026,
\href{https://arxiv.org/abs/2608.01447v3}{arXiv:2608.01447v3}). In the
high-order case the active curvature contributions have the same sign and
reinforce each other; here they have opposite signs and compete.

\subsection{Organization}

Section~\ref{sec:passage} constructs the return map, derives the passage law
and proves the estimates used throughout. Section~\ref{sec:curvature-graph}
proves the bounds of Theorem~\ref{thm:six-main} and their attainment.
Section~\ref{sec:main-saturation} proves
Theorem~\ref{thm:main-saturation-summary}. Section~\ref{sec:actual-bifurcations}
constructs the limit cycles of multiplicity four and proves
Theorem~\ref{thm:actual-quartic-summary}. Section~\ref{sec:sharpness} realizes
the hypotheses by local perturbations and in a harvesting model, and
Section~\ref{sec:concluding} lists open problems.
Appendix~\ref{app:harvesting-fold} constructs the smooth-core fold used in the
harvesting model.
